% =========================================================================================
% Chapter 7 -- System Integration and Prototype
% =========================================================================================
\section{System Integration and Prototype}
\label{sec:integration}

\subsection{End-to-End Architecture}
\label{subsec:integration-arch}

The assembled system is the one shown in Fig.~\ref{fig:architecture}. This section does not
repeat that diagram; it describes what the diagram omits --- where each component actually
runs, what is written to disk, and what happens when a component is not there.

\paragraph{Where each component runs.} The four perception modules load PyTorch checkpoints and
run on a GPU. The reasoning agent loads a 4.9\,GB quantised language model through
\texttt{llama.cpp} and also wants a GPU. These do not fit together in the free Colab session
this project had available: the organ models and an 8-billion-parameter model cannot occupy one
runtime, and attempting it exhausts memory. The agents therefore run at different times, in
different sessions.

That constraint shaped the architecture rather than being worked around. Because the agents
cannot run in one process, they cannot call each other, and the interface between them had to
become a persisted artefact instead of a function call. The consequences turned out to be
advantages, and Section~\ref{subsec:integration-comm} describes them.

\paragraph{What is persisted.} One JSON document per encounter, the \emph{encounter bundle},
holding four sections: the patient's demographics and presenting complaint, the triage agent's
output, one section per organ examined by the ultrasound agent, and the laboratory and vital
values available. The bundle is the complete input to the reasoning agent; nothing else is
consulted, and no image reaches the reasoning layer at any point.

\paragraph{What happens when a component is unavailable.} Every degradation is represented in
the ordinary format rather than as an absence:

\begin{itemize}
  \item an organ with no implemented module returns \texttt{status: "not\_supported"} with an
        empty finding list;
  \item a module whose calibration file is missing runs anyway and reports
        \texttt{confidence\_calibrated: false}, so the reasoning layer discounts it rather than
        treating a raw score as a probability;
  \item a module that fails at inference --- the cardiac module on an image where no ventricle
        is found, for instance --- returns \texttt{status: "failed"} rather than an invented
        measurement;
  \item the language model failing to load, timing out or running out of memory produces a
        withheld differential with the error recorded, while the escalation decision, computed
        before the model is consulted, is returned unchanged.
\end{itemize}

The last case is the one that matters most and it is verified directly rather than assumed: run
with a backend that raises on every call, the pipeline returns the escalation decision with its
triggers intact and withholds only the differential.

\subsection{Agent Communication}
\label{subsec:integration-comm}

The agents communicate through the shared schema of Section~\ref{subsec:output-contract}, which
functions as an interface contract rather than a convention. A perception agent writes through
\texttt{make\_report}, the only exit from a module, and the result is validated before it is
stored: a malformed report is rejected at the boundary rather than discovered downstream.

Writes are atomic. A bundle section is written to a temporary file and moved into place, so a
crash mid-write cannot leave a partially written encounter that still parses as valid JSON. A
truncated bundle read as complete would be worse than a missing one, because nothing downstream
would signal that anything was wrong.

Communicating through stored records rather than direct calls has three consequences worth
stating, because they were what made the rest of this project tractable:

\begin{enumerate}
  \item \textbf{Each agent is independently testable.} The reasoning agent can be exercised
        against a hand-written bundle with no imaging model present at all, which is how all
        281 tests of Section~\ref{subsec:reasoning-safety} run --- in seconds, with no GPU and
        no model weights.
  \item \textbf{Each decision is reproducible after the fact.} The bundle that produced a given
        output is on disk. A decision can be re-examined months later against exactly the inputs
        it was made from, which is a precondition for auditing anything clinical.
  \item \textbf{The agents can be replaced independently.} The reasoning layer does not know
        which model produced a finding, only that the finding arrived through the contract with
        a calibrated confidence and a declared scope.
\end{enumerate}

\subsection{Data Flow}
\label{subsec:integration-flow}

One encounter proceeds through five stages, each producing a named artefact.

\begin{table}[H]
  \centering
  \small
  \caption{Artefacts produced at each stage of one encounter.}
  \label{tab:dataflow}
  \begin{tabularx}{\textwidth}{@{}llX@{}}
    \toprule
    Stage & Artefact & Content \\
    \midrule
    1. Arrival & bundle created & demographics, presenting complaint, vitals, labs \\
    2. Triage & \texttt{triage} section & urgency tier, calibrated confidence, features used \\
    3. Imaging & one section per organ & findings detected, findings screened and not detected,
      confidences, declared scope \\
    4. State & clinical state & the sections merged, plus what is \emph{missing}, the conflicts
      between agents, and a case-quality grade \\
    5. Reasoning & result & escalation decision, evidence list, differential or withheld, the
      abnormal values considered \\
    6. Decision support & support object & severity, critical alerts, recommended
      examinations, routed scenario, therapeutic considerations if a protocol was retrieved \\
    7. Reporting & report + archive & standardised report with a conclusion, written to the
      archive with a content hash \\
    \bottomrule
  \end{tabularx}
\end{table}

Stages 2 and 3 are independent and neither depends on the other, which is why they can run in
either order or in parallel; the triage agent never sees an image and the ultrasound agent
never sees a vital sign. Stage 4 is where the two meet, and it is also where the system first
computes something neither agent could: what is \emph{absent}, and whether the two agents
disagree.

Stage 5 splits internally. The escalation decision is computed from the state alone and is
complete before the language model is consulted. The differential is produced by the model,
checked against the state, and either delivered or withheld. Both leave the system together, so
a withheld differential still arrives with an escalation decision attached.

Stage 6 shares that property and extends it. Severity, alerts, examination recommendations and
scenario routing are computed from the state, not from the model's answer, and are therefore
attached to every result --- including those where the model was never consulted or failed
outright. A report on a case whose differential was withheld still carries a severity and its
alerts, which is the behaviour a clinician needs from a system whose reasoning component has
just declined to answer.

\subsection{Example Clinical Scenarios}
\label{subsec:integration-scenarios}

Three cases are traced end to end below. All three are \textbf{synthetic}, constructed to
exercise specific behaviours rather than drawn from patients; no clinical data was used at any
point in this project. The outputs shown are the actual recorded outputs of the frozen system
described in Section~\ref{subsec:reasoning-ab}, not illustrations of intended behaviour.

\subsubsection{Concordant Case}

Triage and imaging agree, the imaging evidence is strong, and the key laboratory values are
present. The clinical state grades case quality \textsc{good}.

The escalation policy finds no trigger: there is no disagreement between the agents, no
high-risk finding on weak evidence, no key laboratory value absent, and no organ left
unassessed. The route is \emph{direct}.

The model returns \emph{heart failure} at high likelihood, citing four items of evidence by
identifier --- the lung finding and three abnormal laboratory values. The answer passes every
grounding check and is delivered, carrying two warnings: three abnormal vital signs were not
cited anywhere in the reasoning, and the recommendation is phrased as confirming a diagnosis
rather than informing one. Both are reported to the reader alongside the answer, and the three
uncited vitals appear in the evidence list regardless.

This is the case the system is meant to answer directly, and it does.

\subsubsection{Conflicting Case}

The triage agent reports \emph{low} urgency with 0.88 confidence. The cardiac module reports
\emph{severe dysfunction} at 0.74. These cannot both be acted on.

The conflict is detected in the state builder, before any reasoning, by comparing the triage
tier against the severity of the imaging findings. Two escalation triggers fire: the agents
disagree, and a high-risk finding is reported below the confidence at which the policy would
treat it as decisive. The route is \emph{simulation} and the case quality falls to
\textsc{moderate}.

The important property is what the system does \emph{not} do. It does not average the two
assessments, does not defer to whichever agent is more confident, and does not resolve the
disagreement. It has no information with which to decide which agent is wrong, so it makes the
disagreement visible and escalates. The differential is still produced --- \emph{acute coronary
syndrome} --- but it arrives with the conflict attached, and the escalation is what a clinician
sees first.

\subsubsection{Missing Data Case}

A positive lung finding is present, and troponin and lactate were never measured.

Both absent tests appear in the state under \textsc{not measured --- absent, not normal}. They
have no evidence identifier, which means the model cannot cite them for or against any
diagnosis: the constraint is structural rather than an instruction the model is asked to
follow. Escalation fires on positive imaging with key laboratory values absent.

The model returns \emph{pulmonary embolism} at moderate likelihood, and names troponin and
lactate first in \texttt{missing\_information} --- the tests that would most change the
assessment. That list is what determines what gets ordered next, which is why an entry that is
merely plausible is not good enough there and why the validator checks it against what the
record actually lacks.

\subsection{User Interface}
\label{subsec:integration-ui}

A browser application was built with Streamlit. Its design follows from a measurement rather
than from a mock-up: Section~\ref{subsec:integration-execution} shows that everything except
the language model completes in under two seconds, which means the safety layer is fast enough
to be \emph{interactive} while the differential is not.

The interface is built around that asymmetry. It is organised as a clinician would work rather
than as the pipeline is built: a four-step intake --- patient, clinical measurements,
point-of-care ultrasound, assessment --- followed by screens for the alerts, the reasoning
chain, the retrieved sources, an encounter timeline and the generated report. Rather than
replaying fixed cases, it lets the user enter a patient: the presenting complaint, the triage
tier, which organs were requested, which were actually scanned, and for every vital and
laboratory value a checkbox that makes it \emph{never measured} rather than normal. The
clinical state, the conflicts, the list of what is absent, the evidence identifiers, the
severity, the alerts and the escalation decision are recomputed on every change, with no
language model and no GPU involved.

The point-of-care ultrasound step offers the three organs for which modules were trained ---
lung, cardiac and gallbladder --- and no others. The output contract admits vascular and FAST
as organs, but no module was built for either, and an examination tab with nothing behind it
would advertise a capability that does not exist. On a screen whose purpose is to distinguish
\emph{not detected} from \emph{never examined}, that would be the wrong claim to make.

Four things about it are worth reporting.

\paragraph{The module reads the image.} The clinician uploads a study and the module analyses it
in the browser, on CPU, in a few seconds. What appears on screen is the module's own report ---
which findings crossed their tuned operating points, which were screened for and not seen, the
thresholds used and their provenance, whether the confidence is calibrated and what ceiling the
calibrator imposes --- and none of it is typed in. This is the step that makes the system
multimodal in operation rather than only in architecture.

All three organs do this, and the interface differs between them because the modules differ. The
lung takes one image and reports four independent findings. The gallbladder takes one image and
reports exactly one of five classes, since it is single-label. The cardiac module takes
\emph{two} images, end-diastole and end-systole, because an ejection fraction is a comparison
between them and no single still can produce one; offered one frame it says so rather than
returning a number. Where a module cannot run at all, the interface asks instead what it
reported and states which of the two obstacles applies --- the checkpoint being absent, or
inference for that organ not having been extracted --- because those have different remedies and
a single message that conflated them would be a false explanation of a true refusal.

\paragraph{Absence is visible as absence.} Unticking a laboratory value moves it out of the
evidence list and into a panel headed \emph{never measured --- no identifier exists}. The
consequence is shown rather than described: with no identifier, that value cannot be cited for
or against any diagnosis, and the escalation decision changes in front of the user.

\paragraph{The failure path is executed, not illustrated.} A control labelled \emph{break the
model} calls \texttt{reason} with a backend that raises on every invocation. What appears is
the system's actual output: the differential withheld with the backend error recorded, and the
escalation decision unchanged beside it. This is the central architectural claim of the project,
and the interface demonstrates it by performing it.

\paragraph{Live and recorded are labelled distinctly, and the label is enforced.} The
differential requires a 4.9\,GB model on a GPU, so for the five benchmark encounters the
interface displays the answers recorded from the frozen run of
Section~\ref{subsec:reasoning-ab}, while everything computed in the browser is marked as
computed there. An interface that blurred the two would misrepresent what the machine in front
of the examiner is doing.

The label is not merely written on the screen. A recorded differential belongs to one specific
record, so the interface attaches it only when the encounter on screen is still \emph{that}
record: if the clinician edits any field, or an uploaded image has been analysed, the recording
is dropped and the screen says that no differential was generated. An untouched benchmark
encounter is built by the same function the regression harness uses, rather than by
reconstructing an equivalent-looking bundle in the interface --- the reliability text alone
would differ, and that text reaches the prompt.

This guard caught a defect of its own during development, which is the reason for describing
it. The wizard's benchmark presets had been written out beside the canonical scenarios rather
than derived from them, and had drifted: a different age, a sex, one missing blood pressure. A
recorded differential was therefore about to appear next to a record it had not been computed
from. The presets are now derived from the scenarios, and the alert counts the interface
reports agree exactly with those the regression harness reports for the same encounters.

\subsubsection{Where the visual design and the system disagreed}

The application's visual design --- its palette, typography, layout and screen structure ---
was produced separately from the system and adopted as delivered. Two pieces of its
\emph{placeholder content} could not be reproduced literally, and the reason is worth recording
because it is the same principle the rest of the architecture rests on.

The first was a roster of patients: named individuals, a count of assessments completed today,
a table of recent cases with clinical statuses. These are ordinary mock data in a design, and
in a clinical interface they are fabricated records. The same tables and cards are rendered in
the same style from the project's own benchmark encounters, whose severities and alert counts
are computed rather than written.

The second was a conversational assistant supplied with a dictionary of pre-written replies. A
chat that answers a typed clinical question with a fluent paragraph implies a conversational
model reasoning about the patient, and this system has none: its language model emits one
structured, validated differential and nothing else. The screen is built as designed, but every
answer is assembled from the computed assessment --- the escalation triggers, the alerts, the
absent tests, the evidence identifiers actually cited --- and a free-typed question is answered
by saying plainly that only the computed assessment can be answered from.

The general point is that a demonstration interface is part of the safety argument, not a
presentation layer above it. A screen that overstates what the system did is a failure of the
same kind as a differential citing a test that was never performed; it is simply committed by
the interface rather than by the model.

Alongside the application, three lower-level interfaces remain, each honest about its fidelity:

\begin{itemize}
  \item \textbf{Notebooks.} The perception modules are trained in Jupyter notebooks on Colab,
        one per organ, and all reported training results were produced through them. Inference
        no longer depends on a notebook for any organ: it was extracted into an importable
        module, which is what the application calls and what the test suite exercises
        (Section~\ref{subsubsec:agent-extraction}).
  \item \textbf{A command-line entry point.} \texttt{run\_case} executes a complete encounter
        end to end, with switches for the retrieval backend, the number of correction rounds
        and whether to consult the model at all. Its \texttt{-{}-dry-run} mode prints the
        clinical state, the escalation decision and the exact prompt without loading the
        language model, which is how the safety layer is inspected without a GPU.
  \item \textbf{A rendered clinical summary.} The system produces a deterministic text
        rendering of the encounter --- state, escalation, evidence considered, differential or
        withheld. The application's report screen displays it and offers it for export.
\end{itemize}

The rendering underlies what the application displays, and it was written deterministically
for that reason: the same bundle always produces byte-identical output, so what a clinician
reads is reproducible and auditable rather than regenerated each time.

What the application is \emph{not} is also worth stating. It has no authentication, no patient
record, no persistence between sessions and no connection to any hospital system. It runs the
prototype's own synthetic encounters on a laptop, and it would need all of the above before it
could be shown a real patient --- which the limitations of Section~
ef{subsec:results-limitations}
forbid in any case.

\subsection{End-to-End Execution}
\label{subsec:integration-execution}

Table~\ref{tab:runtime} gives measured times for a complete encounter on the hardware this
project had available: a Colab T4 GPU.

\begin{table}[H]
  \centering
  \small
  \caption{Measured runtime for one encounter, Colab T4.}
  \label{tab:runtime}
  \begin{tabular}{@{}lr@{}}
    \toprule
    Stage & Time \\
    \midrule
    Perception model loading (three organs) & $\sim$10\,s \\
    Per-organ inference & $<$1\,s \\
    State construction and escalation & $<$0.1\,s \\
    Retrieval over the corpus & $<$0.1\,s \\
    Language model loading (4.9\,GB) & 14--21\,s \\
    Reasoning, including one correction round & 33--105\,s \\
    \bottomrule
  \end{tabular}
\end{table}

Two features of this profile constrain bedside use.

The first is that everything except the language model is fast. Perception, state construction,
conflict detection and the entire escalation decision complete in under two seconds. A
deployment that offered only the escalation decision --- which is the safety-critical output,
and the one that does not depend on the model --- would be effectively instantaneous.

The second is that the reasoning step dominates and is variable, from 33 to 105 seconds
depending on the case and whether a correction round was needed. On CPU the same workload takes
approximately 23 minutes per case, which rules out CPU-only deployment entirely.

The practical implication is that the two outputs have different latency profiles and should
probably be delivered differently. The escalation decision is available immediately and is what
a clinician needs first; the differential and its reasoning arrive up to two minutes later and
are what they would read while waiting. The architecture already separates these two paths for
safety reasons, and the timing measurements suggest the same separation is right for
interaction as well.

\subsection{A Controlled Path for Future Model Updates}
\label{subsec:mlops}

The datasets behind the perception models do not cover the whole clinical scope this system is
intended for. That is ordinary for medical imaging work at this scale, and it has a
consequence for the architecture rather than only for the results: a system whose models are
expected to improve later needs a defined mechanism for improving them, or the first
contribution of new validated data becomes an ad-hoc retraining whose provenance nobody can
reconstruct.

What is implemented is deliberately small. New validated data updates a dataset manifest; the
manifest is content-hashed into a dataset version; a training command consumes that version
and writes its own metrics; the run is registered as a \emph{candidate}; and an automated gate
decides whether the candidate may be offered for approval. A person then promotes it, or does
not. Figure~\ref{fig:mlops} shows the cycle.

\begin{figure}[H]
  \centering
  \begin{tikzpicture}[
      node distance = 6mm,
      font = \small,
      box/.style    = {draw, rounded corners=2pt, align=center, inner sep=4pt,
                       minimum height=7mm, text width=42mm},
      data/.style   = {box, fill=black!4},
      run/.style    = {box, fill=black!8},
      gateb/.style  = {box, fill=black!12, very thick, text width=46mm},
      human/.style  = {box, fill=black!12, very thick, text width=42mm},
      term/.style   = {box, text width=26mm},
      side/.style   = {box, dashed, fill=none, text width=30mm},
      ar/.style     = {-{Latex[length=2mm]}, thick},
      lbl/.style    = {font=\scriptsize\itshape, inner sep=1pt},
    ]

    \node[data] (new) {new validated data};
    \node[data, below=of new] (ver) {dataset manifest\\ \texttt{lung@11075e6a4296}};
    \node[run,  below=of ver] (train) {reproducible training};
    \node[run,  below=of train] (eval) {evaluation\\ \emph{fixed protocol}};
    \node[run,  below=of eval] (reg) {register \textbf{candidate}};
    \node[gateb, below=of reg] (gate) {\textbf{promotion gate}\\
      headline $\cdot$ guarded recall\\ floors $\cdot$ calibration\\
      same eval set $\cdot$ dataset recorded};

    \draw[ar] (new) -- (ver);
    \draw[ar] (ver) -- (train);
    \draw[ar] (train) -- (eval);
    \draw[ar] (eval) -- (reg);
    \draw[ar] (reg) -- (gate);

    % MLflow sits beside the flow: it records, it does not decide.
    \node[side, right=10mm of reg] (mlf) {MLflow\\ \scriptsize experiment tracking};
    \draw[ar, dashed] (reg) -- (mlf);

    % The two outcomes.
    \node[term, below left=10mm and 2mm of gate] (rej) {\textbf{rejected}};
    \node[human, below right=10mm and 2mm of gate] (appr) {\textbf{human approval}\\
      \scriptsize named approver, recorded};
    \draw[ar] (gate.south) -| node[lbl, pos=0.25, above] {fail} (rej.north);
    \draw[ar] (gate.south) -| node[lbl, pos=0.25, above] {pass} (appr.north);

    \node[term, below=of appr] (prod) {\textbf{production}};
    \draw[ar] (appr) -- (prod);

  \end{tikzpicture}
  \caption{The model-update lifecycle. The gate decides only whether a candidate may be
  \emph{offered}; no path leads from contributed data to a deployed model without a named
  person approving it. MLflow records runs and takes no part in the decision.}
  \label{fig:mlops}
\end{figure}

\subsubsection{Dataset Versioning}

A dataset version is the SHA-256 of the manifest bytes, truncated for readability
(\texttt{lung@11075e6a4296}). The manifest is the list of rows actually used, so the hash
changes exactly when the training data changes, and every registered model records the version
it was trained on.

DVC was considered and not adopted. The media is tens of gigabytes, is already untracked, and
is not redistributable under the source licences; DVC without a storage backend would have
added a tool without adding reproducibility. Hashing the manifest gives the property that
matters --- which rows produced which model --- and \texttt{dvc add} on the manifests slots in
underneath the same design if a suitable backend becomes available.

\subsubsection{Experiment Tracking and the Model Registry}

Runs are logged to MLflow, backed by SQLite so that no server is required. The registry itself
is a JSON file recording, per model version: the stage (\emph{candidate}, \emph{production},
\emph{archived}), the dataset version, the evaluation set, the metrics, the artifact path, the
training configuration, and --- once promoted --- who approved it and when.

MLflow is optional by construction. A tracking backend that is missing or misconfigured is
reported and stepped over rather than raised, because a completed training run must not be
lost to a logging failure. Version names are never reused: the counter only advances, retired
names are carried across re-seeds of the registry, and an explicit name that collides is
refused. This was added after MLflow, which is an append-only log and was right to keep both,
came to hold two different runs called \texttt{v3}.

\subsubsection{The Promotion Gate}

\emph{Is the new model better?} is the wrong question for an emergency system, because a model
can gain aggregate accuracy by classifying easy cases better while missing more of the patients
who are actually ill. The gate therefore checks what a single headline number hides:

\begin{itemize}
  \item the headline metric, which may not regress;
  \item \textbf{guarded metrics} --- recall on the findings where a miss is dangerous ---
        which may fall by at most a stated tolerance;
  \item absolute floors;
  \item calibration, because the reasoning layer escalates on low confidence, so a model whose
        stated confidence stops meaning anything breaks the escalation logic while appearing
        to classify better;
  \item that the candidate was scored on the \textbf{same evaluation set} as the incumbent,
        since two numbers from two protocols are two facts rather than an improvement;
  \item that the dataset version is recorded, because a model that cannot be reproduced cannot
        be withdrawn and rebuilt when something is found wrong.
\end{itemize}

A missing guard metric fails rather than passes. A candidate that does not report its
dangerous-miss rate has not shown that it did not regress, and silence must not read as
success --- the same rule the clinical state applies to an unmeasured laboratory value.

Two commands ask two different questions. \texttt{gate} asks whether a candidate may replace
what is deployed. \texttt{compare} asks whether one run improved on another, and changes
nothing; it exists so that establishing a pipeline does not require promoting a model to
production in order to measure it, which would make the registry state something false about
what a clinician is seeing.

\subsubsection{Human Approval}

Training and gating stop at \emph{candidate}. Promotion requires a named approver and records
the name, the reason and the time; without one it is refused. A test asserts that the training
command cannot write a production stage, call the promotion path, or set an approver, so the
dangerous version of this design --- contributed data reaching deployed models unattended ---
is not reachable by accident rather than merely discouraged by convention.

The protocol above is frozen in \texttt{models/mlops\_protocol.json} and compared against the
enforcing code by a test. Changing a threshold requires re-freezing with a recorded reason,
because a threshold revised after seeing the result it decides is not a threshold; it is a
description of the result.
